\documentclass[../../main.tex]{subfiles}

\begin{document}

\chapter{Arsitektur Mikroprosesor 8086: BIU dan EU}
\label{chap:arsitektur-8086}

% =========================================================
% OUTCOME PEMBELAJARAN KHUSUS BAB
% =========================================================

\begin{subcpmkbox}
\vspace{1.1em}
\SubCPMK{1.2} --- Menjelaskan pembagian \textit{Bus Interface Unit}
dan \textit{Execution Unit}, peran bus, \textit{queue} instruksi, dan
ALU, serta urutan pengambilan dan eksekusi instruksi pada 8086.

\vspace{0.4em}
\textbf{CPMK yang didukung:} \CPMK{01} \quad\textbullet\quad
\textbf{CPL yang didukung:} \CPL{02}
\end{subcpmkbox}

\begin{tujuanbox}
\vspace{1.1em}
Setelah mempelajari bab ini, mahasiswa mampu:
\begin{enumerate}
    \item menjelaskan kedudukan mikroprosesor 8086 dan alasan
          pembagian internalnya menjadi dua unit;
    \item menjelaskan tugas \eng{Bus Interface Unit} (BIU) beserta
          komponen bus, \eng{queue} instruksi, register segmen, dan
          \eng{Instruction Pointer};
    \item menjelaskan tugas \eng{Execution Unit} (EU) beserta ALU,
          \eng{flag}, \eng{instruction decoder and control}, dan
          register yang terlibat;
    \item menghitung alamat fisik dari pasangan \eng{segment:offset}
          sebagai bagian dari kerja BIU;
    \item menelusuri urutan pengambilan, penerjemahan, dan eksekusi
          satu instruksi;
    \item membedakan instruksi yang seluruhnya bekerja di register
          dari instruksi yang memerlukan akses memori melalui BIU.
\end{enumerate}
\end{tujuanbox}

\begin{materibox}
\vspace{1.1em}
\begin{enumerate}
    \item Kedudukan 8086 dan pembagian unit internal.
    \item \eng{Bus Interface Unit}: tugas dan komponen.
    \item Bus alamat, bus data, dan bus kontrol.
    \item \eng{Instruction queue} enam byte dan pengisiannya.
    \item \eng{Execution Unit}: ALU, \eng{flag}, \eng{operand}, dan
          decoder.
    \item Urutan \eng{fetch}--\eng{decode}--\eng{execute}.
\end{enumerate}
\end{materibox}

\subfile{section-01}
\subfile{section-02}
\subfile{section-03}

\section{Contoh Program}
\label{sec:contoh-bab4}

Dua contoh berikut dipakai untuk menguji pemahaman tentang siklus
instruksi. Contoh pertama memakai program sesingkat mungkin agar
seluruh langkah dapat ditelusuri, sedangkan contoh kedua
membandingkan dua cara penjumlahan yang berbeda biayanya.

\subsection{Contoh 1: Menelusuri Siklus Instruksi sebuah Program}

Kode~\ref{lst:siklus-dasar} adalah program utuh yang berisi satu
instruksi utama, yaitu \lstasm|MOV AX, 1234h|. Program ini sengaja
dibuat pendek agar setiap langkah pada Emu8086 dapat diamati tanpa
kehilangan jejak.

\lstinputlisting[
    language=MASM,
    caption={Program \file{siklus-dasar.asm} untuk penelusuran siklus
    instruksi.},
    label={lst:siklus-dasar}
]{\CodeDir/siklus-dasar.asm}

Setelah program dirakit, hasil \eng{disassembly} pada Emu8086
memperlihatkan setiap instruksi beserta panjangnya dalam byte.
Panjang tersebut penting karena langsung menentukan berapa kali
\reg{IP} bertambah dan berapa banyak byte yang masuk ke
\eng{instruction queue}. Tabel~\ref{tab:trace-program} mencatat
penelusuran program tersebut dari titik masuk hingga program selesai.

\begin{table}[H]
    \centering
    \caption{Hasil penelusuran program \file{siklus-dasar.asm}
    dengan \eng{instruction queue} berisi maksimum enam byte.}
    \label{tab:trace-program}
    \small
    \begin{tabular}{@{}clp{1.5cm}p{2.3cm}p{2.0cm}p{3.6cm}@{}}
        \toprule
        \textbf{Langkah} & \textbf{\reg{CS}:\reg{IP}} &
        \textbf{Panjang} & \textbf{Instruksi} &
        \textbf{Antrean} & \textbf{Keadaan setelahnya} \\
        \midrule
        1 & \hex{1000}:\hex{0000} & 3 byte & \code{MOV AX, @DATA} &
        \hex{B8 00 10} & \reg{AX} $=$ \hex{1000} \\
        \addlinespace
        2 & \hex{1000}:\hex{0003} & 2 byte & \code{MOV DS, AX} &
        \hex{8E D8} & \reg{DS} $=$ \hex{1000} \\
        \addlinespace
        3 & \hex{1000}:\hex{0005} & 3 byte & \code{MOV AX, 1234h} &
        \hex{B8 34 12} & \reg{AX} $=$ \hex{1234} \\
        \addlinespace
        4 & \hex{1000}:\hex{0008} & 2 byte & \code{MOV AH, 4Ch} &
        \hex{B4 4C} & \reg{AH} $=$ \hex{4C} \\
        \addlinespace
        5 & \hex{1000}:\hex{000A} & 2 byte & \code{MOV AL, 00h} &
        \hex{B0 00} & \reg{AL} $=$ \hex{00} \\
        \addlinespace
        6 & \hex{1000}:\hex{000C} & 2 byte & \code{INT 21h} &
        \hex{CD 21} & Program berakhir \\
        \bottomrule
    \end{tabular}
\end{table}

Tiga hal dapat disimpulkan dari Tabel~\ref{tab:trace-program}.

\begin{enumerate}
    \item \textbf{\reg{IP} bertambah menurut panjang instruksi.}
          Instruksi tiga byte menambah \reg{IP} sebanyak tiga,
          sedangkan instruksi dua byte menambahnya sebanyak dua.

    \item \textbf{Antrean menampung lebih dari satu instruksi.}
          Pada langkah 4 dan 5, kedua instruksi hanya berukuran dua
          byte sehingga keduanya dapat berada dalam antrean enam
          byte sekaligus.

    \item \textbf{Instruksi yang mengakses memori menambah siklus
          bus.} Instruksi yang membaca variabel pada segmen data
          memicu siklus bus tambahan yang tidak dimiliki instruksi
          register-ke-register.
\end{enumerate}

\subsection{Contoh 2: Dua Cara Menjumlahkan Dua Bilangan}

Contoh kedua membandingkan biaya dua pendekatan penjumlahan. Pada
cara pertama, kedua operand sudah berada di register. Pada cara
kedua, operand dibaca dari memori sehingga BIU harus menjalankan
siklus bus. Program lengkapnya ditunjukkan pada
Kode~\ref{lst:jumlah-dua-cara}.

\lstinputlisting[
    language=MASM,
    caption={Program \file{jumlah-dua-cara.asm}: penjumlahan
    melalui register dan melalui memori.},
    label={lst:jumlah-dua-cara}
]{\CodeDir/jumlah-dua-cara.asm}

Kedua cara menghasilkan nilai yang sama, yaitu \hex{1234} $+$
\hex{5678} $=$ \hex{68AC}. Meskipun demikian, jumlah siklus yang
dibutuhkan berbeda. Tabel~\ref{tab:perbandingan-siklus} membandingkan
perkiraan jumlah siklus \eng{clock} untuk setiap instruksi
berdasarkan datasheet 8086 \cite{intel1979}.

\begin{table}[H]
    \centering
    \caption{Perkiraan jumlah siklus \eng{clock} kedua cara
    penjumlahan.}
    \label{tab:perbandingan-siklus}
    \small
    \begin{tabular}{@{}p{3.6cm}cp{4.4cm}p{3.4cm}@{}}
        \toprule
        \textbf{Instruksi} & \textbf{Siklus} & \textbf{Jenis operand} &
        \textbf{Akses bus} \\
        \midrule
        \multicolumn{4}{@{}l}{\textbf{Cara 1 --- operand di
        register}} \\
        \addlinespace
        \code{MOV AX, 1234h} & 4 & nilai segera &
        hanya pengambilan instruksi \\
        \addlinespace
        \code{MOV BX, 5678h} & 4 & nilai segera &
        hanya pengambilan instruksi \\
        \addlinespace
        \code{ADD AX, BX} & 3 & register--register &
        tidak ada \\
        \addlinespace
        \textbf{Jumlah cara 1} & \textbf{11} & & \\
        \midrule
        \multicolumn{4}{@{}l}{\textbf{Cara 2 --- operand di memori}} \\
        \addlinespace
        \code{MOV AX, bil1} & 14 & memori langsung &
        satu baca memori \\
        \addlinespace
        \code{ADD AX, bil2} & 15 & register--memori &
        satu baca memori \\
        \addlinespace
        \textbf{Jumlah cara 2} & \textbf{29} & & \\
        \bottomrule
    \end{tabular}
\end{table}

Perbandingan pada Tabel~\ref{tab:perbandingan-siklus} menunjukkan
bahwa cara kedua memerlukan waktu hampir tiga kali lebih lama untuk
pekerjaan yang sama. Penyebabnya bukan pada operasi penjumlahannya,
melainkan pada dua siklus bus tambahan yang harus dijalankan BIU.
Karena BIU juga sibuk mengambil instruksi berikutnya, siklus bus
yang panjang dapat mengosongkan antrean dan meniadakan keuntungan
\eng{pipeline}.

\begin{infobox}
Angka pada Tabel~\ref{tab:perbandingan-siklus} adalah nilai khas
untuk 8086 dengan pengalamatan langsung. Nilai sebenarnya bergantung
pada mode pengalamatan yang dipakai, karena setiap mode memiliki
biaya perhitungan alamat (\eng{effective address}) yang berbeda.
Cara penerjemahan instruksi menjadi kode mesin dibahas pada
Bab~\ref{chap:format-instruksi}.
\end{infobox}

\begin{session}[caption={Keluaran pemeriksaan isi memori
\file{jumlah-dua-cara.asm} pada Emu8086.}, label={lst:keluaran-hasil}]
1000:0008  hasil1  DW  68AC
1000:000A  hasil2  DW  68AC
\end{session}

Kedua penampung berisi nilai yang sama, tetapi waktu tempuhnya
berbeda. Inilah pelajaran utama bab ini: memahami arsitektur berarti
memahami \emph{mengapa} dua program yang hasilnya sama bisa memiliki
biaya yang berbeda.

\section{Praktikum}
\label{sec:praktikum-bab4}

\begin{praktikumbox}
\vspace{1.1em}
\textbf{Judul:} Diagram blok 8086 dan penelusuran siklus instruksi.\\[0.3em]
\textbf{Sub-CPMK:} 1.2 \quad\textbullet\quad
\textbf{Alokasi waktu:} 100 menit \quad\textbullet\quad
\textbf{Bentuk:} individu
\end{praktikumbox}

\subsection{Tujuan Praktikum}

Setelah menyelesaikan praktikum ini, mahasiswa mampu menggambar
diagram blok BIU/EU, menjelaskan pembagian tugas kedua unit, serta
menelusuri perubahan \reg{IP}, \reg{AX}, dan isi \eng{queue} pada
setiap langkah eksekusi.

\subsection{Alat dan Bahan}

\begin{itemize}
    \item Komputer dengan emu8086 terpasang.
    \item DOSBox dengan TASM, TLINK, dan DEBUG.
    \item Berkas \file{siklus-dasar.asm} dari \file{code/asm/}.
    \item Kertas kosong, alat tulis, dan lembar jurnal praktikum.
\end{itemize}

\subsection{Langkah Kerja}

\begin{enumerate}
    \item \textbf{Menggambar diagram blok dari ingatan.}
    \begin{enumerate}
        \item Tutup halaman materi, lalu gambarkan kotak EU dan BIU
              beserta bus internal dan bus eksternal.
        \item Tuliskan komponen yang berada di dalam EU dan BIU,
              termasuk \eng{instruction queue} enam byte.
        \item Buka kembali Gambar~\ref{fig:blok-8086} dan tandai
              bagian yang salah letak atau terlewat.
    \end{enumerate}

    \item \textbf{Menelusuri program dengan emu8086.}
    \begin{enumerate}
        \item Buka \file{siklus-dasar.asm}, lalu pilih
              \file{Compile and Emulate}.
        \item Catat nilai awal \reg{CS}, \reg{IP}, \reg{AX}, dan
              \reg{DS}.
        \item Tekan \file{Single Step} satu kali. Catat perubahan
              \reg{IP}, \reg{AX}, dan isi \eng{queue}.
        \item Ulangi sampai instruksi \code{INT 21h}. Hentikan
              penelusuran sebelum layanan DOS dijalankan.
        \item Susun hasil pengamatan dalam tabel seperti
              Tabel~\ref{tab:trace-program}.
    \end{enumerate}

    \item \textbf{Menelusuri dengan DEBUG.}
    \begin{enumerate}
        \item Dari DOSBox, rakit program dengan
              \file{tasm siklus-dasar.asm}, lalu link dengan
              \file{tlink siklus-dasar.obj}.
        \item Jalankan \file{debug siklus-dasar.exe}.
        \item Ketik \file{r} untuk menampilkan register.
        \item Ketik \file{t} untuk menjalankan satu instruksi, lalu
              catat perubahan \reg{IP} dan \reg{AX}.
        \item Bandingkan hasil DEBUG dengan hasil emu8086.
    \end{enumerate}

    \item \textbf{Menentukan unit yang bekerja.}
    \begin{enumerate}
        \item Untuk setiap instruksi pada tabel, tandai apakah
              pengambilan instruksi dikerjakan BIU, eksekusi
              dikerjakan EU, atau keduanya.
        \item Jelaskan mengapa \code{MOV AX, 1234h} tidak memerlukan
              siklus bus data tambahan.
    \end{enumerate}

    \item \textbf{Modifikasi program.}
    \begin{enumerate}
        \item Tambahkan instruksi \code{MOV BX, AX} setelah
              \code{MOV AX, 1234h}.
        \item Rakit ulang, lalu telusuri kembali.
        \item Catat perubahan \reg{BX} dan panjang instruksi yang
              ditambahkan.
    \end{enumerate}

    \item \textbf{Menghitung alamat fisik.}
    \begin{enumerate}
        \item Untuk setiap baris tabel, hitung alamat fisik
              \reg{CS}:\reg{IP} dengan rumus
              $(\reg{CS} \times 16) + \reg{IP}$.
        \item Tuliskan hasilnya dalam kolom tambahan.
    \end{enumerate}
\end{enumerate}

\subsection{Hasil yang Dikumpulkan}

\begin{enumerate}
    \item Gambar diagram blok BIU/EU beserta catatan koreksi.
    \item Tabel pengamatan \reg{CS}, \reg{IP}, \reg{AX}, \reg{DS},
          dan isi \eng{queue} untuk setiap langkah.
    \item Tabel alamat fisik yang dihitung dari pasangan
          \reg{CS}:\reg{IP}.
    \item Berkas \file{siklus-dasar.asm} versi modifikasi.
    \item Analisis singkat tentang perbedaan pekerjaan BIU dan EU.
\end{enumerate}

\clearpage
\begin{asesmenbox}
\vspace{1.1em}
\textbf{Asesmen praktikum 2 (bobot 5\% dari komponen praktikum).}

\begin{tabular}{@{}p{6.0cm}cp{4.8cm}@{}}
    \toprule
    \textbf{Indikator} & \textbf{Bobot} & \textbf{Kriteria} \\
    \midrule
    Ketepatan diagram blok BIU/EU & 30\% &
    Seluruh komponen utama berada pada unit yang benar \\
    \addlinespace
    Ketepatan tabel penelusuran & 30\% &
    Nilai \reg{IP}, \reg{AX}, dan isi antrean tercatat benar \\
    \addlinespace
    Ketepatan perhitungan alamat fisik & 25\% &
    Seluruh alamat dihitung dengan rumus yang benar \\
    \addlinespace
    Analisis pembagian tugas & 15\% &
    Perbedaan kerja BIU dan EU dijelaskan dengan tepat \\
    \bottomrule
\end{tabular}
\end{asesmenbox}

\section{Latihan}
\label{sec:latihan-bab4}

\begin{enumerate}
    \item \textbf{Konseptual.} Jelaskan perbedaan antara arsitektur
          internal dan arsitektur eksternal sebuah mikroprosesor.
          Sebutkan dua contoh hal yang termasuk arsitektur internal
          dan dua contoh yang termasuk arsitektur eksternal pada
          8086.

    \item \textbf{Konseptual.} Mengapa 8086 dibagi menjadi dua unit,
          yaitu \eng{Bus Interface Unit} dan \eng{Execution Unit}?
          Jelaskan keuntungan pembagian ini dibandingkan prosesor
          yang hanya memiliki satu unit pengolah.

    \item \textbf{Perhitungan alamat.} Diketahui \reg{CS} $=$
          \hex{2A10} dan \reg{IP} $=$ \hex{00F3}. Hitung alamat fisik
          instruksi berikutnya dengan memakai
          Rumus~\ref{eq:alamat-fisik}. Tunjukkan langkah
          perhitungannya, bukan hanya hasil akhirnya.

    \item \textbf{Analisis \eng{pipeline}.} Jelaskan mengapa
          instruksi lompatan seperti \asm{JMP} menghapus keuntungan
          \eng{instruction queue}. Kaitkan jawaban Anda dengan
          istilah \eng{queue flush} dan sebutkan akibatnya pada
          program yang banyak memuat lompatan.

    \item \textbf{Analisis kode.} Perhatikan potongan program
          berikut.

\begin{lstlisting}[language=MASM, numbers=none]
    MOV AX, 0001h
    MOV BX, 0002h
    ADD AX, BX
    MOV CX, AX
\end{lstlisting}

          \begin{enumerate}
              \item Berapa byte total keempat instruksi tersebut?
              \item Berapa nilai \reg{IP} setelah keempat instruksi
                    selesai bila \reg{IP} mula-mula bernilai
                    \hex{0100}?
              \item Pada instruksi manakah tahap \eng{store}
                    benar-benar memerlukan siklus bus? Jelaskan.
          \end{enumerate}

    \item \textbf{Analisis bus.} Lengkapi pernyataan berikut dan
          jelaskan alasan setiap jawaban.
          \begin{enumerate}
              \item Bus yang bersifat satu arah adalah bus \dots
              \item Bus yang lebarnya menentukan kapasitas memori
                    maksimum adalah bus \dots
              \item Sinyal yang menandai bahwa pin \reg{AD0}--\reg{AD15}
                    sedang membawa alamat adalah \dots
              \item Sinyal yang menentukan sasaran akses memori atau
                    I/O adalah \dots
          \end{enumerate}

    \item \textbf{Penulisan program.} Tulislah program lengkap yang
          menyalin nilai \hex{ABCD} dari register \reg{AX} ke
          register \reg{BX}, lalu menambahkan \hex{0001} pada
          \reg{BX}. Jelaskan pada langkah mana tahap \eng{fetch},
          \eng{decode}, \eng{execute}, dan \eng{store} terjadi untuk
          instruksi terakhir.

    \item \textbf{Analisis kesalahan.} Seorang mahasiswa mengamati
          bahwa \reg{IP} pada Emu8086 bertambah tiga setelah
          menjalankan instruksi \lstasm|MOV AX, 1234h|, tetapi hanya
          bertambah satu setelah menjalankan instruksi
          \lstasm|INC AX|. Mahasiswa tersebut menyimpulkan bahwa
          programnya salah. Tentukan apakah kesimpulan itu benar,
          lalu jelaskan alasan sebenarnya.
\end{enumerate}

\begin{refleksibox}
Jawab pertanyaan berikut pada jurnal praktikum Anda.
\begin{enumerate}
    \item Bagian mana dari bab ini yang paling sulit dibayangkan
          secara visual: diagram blok, \eng{pipeline}, atau bus
          eksternal? Apa yang Anda lakukan untuk mengatasinya?
    \item Setelah menelusuri satu instruksi dengan \eng{Single Step},
          apakah pemahaman Anda tentang cara kerja \reg{IP}
          berubah? Jelaskan bagian mana yang paling mengejutkan.
    \item Konsep apa yang masih perlu Anda jelaskan ulang sebelum
          masuk ke Bab~\ref{chap:register}?
\end{enumerate}
\end{refleksibox}

\section{Rangkuman}
\label{sec:rangkuman-bab4}

Pada bab ini telah dibahas:

\begin{itemize}
    \item Mikroprosesor berkembang dari Intel 4004 (4 bit) hingga
          8086 (16 bit) yang menjadi dasar keluarga x86. Varian 8088
          dengan bus luar 8 bit dipakai pada IBM PC.

    \item Arsitektur dipandang dari dua sisi: \term{internal}
          (register, ALU, dan unit di dalam chip) dan
          \term{eksternal} (pin, bus, dan sinyal kendali).

    \item 8086 terbagi menjadi \eng{Bus Interface Unit} (BIU) dan
          \eng{Execution Unit} (EU). BIU memegang seluruh lalu lintas
          bus luar; EU hanya menjalankan instruksi dan tidak
          terhubung langsung ke pin luar.

    \item BIU memuat register segmen (\reg{CS}, \reg{DS}, \reg{SS},
          \reg{ES}), \reg{IP}, penghitung alamat 20-bit, dan
          \eng{instruction queue} enam byte.

    \item EU memuat \eng{ALU}, register umum, register penunjuk dan
          indeks, serta register \eng{flag}.

    \item Rumus alamat fisik adalah (segmen $\times$ 16) $+$ offset,
          sehingga 8086 dapat menjangkau 1 MB memori meskipun
          register-nya hanya 16 bit.

    \item \eng{Instruction queue} memungkinkan \eng{pipelining}:
          pengambilan instruksi berikutnya berimpit dengan eksekusi
          instruksi sekarang. Instruksi lompatan mengosongkan
          antrean (\eng{queue flush}) dan menghapus keuntungan ini.

    \item Tiga kelompok bus eksternal adalah bus alamat 20-bit
          (satu arah), bus data 16-bit (dua arah), dan bus kontrol.
          Pin \reg{AD0}--\reg{AD15} dimultipleks sehingga alamat
          harus ditangkap dengan bantuan sinyal \reg{ALE}.

    \item Siklus instruksi terdiri atas \eng{fetch}, \eng{decode},
          \eng{execute}, dan \eng{store}. Tahap \eng{store} hanya
          memerlukan siklus bus bila tujuannya berada di memori.
\end{itemize}

\begin{checklistbox}
Tandai setiap butir yang sudah Anda kuasai.

\begin{itemize}
    \item[\checkbox] Saya dapat menjelaskan perbedaan arsitektur
          internal dan eksternal sebuah mikroprosesor.
    \item[\checkbox] Saya dapat menggambar diagram blok BIU/EU
          beserta isi register masing-masing unit.
    \item[\checkbox] Saya dapat menjelaskan peran
          \eng{instruction queue} enam byte dan \eng{pipelining}.
    \item[\checkbox] Saya dapat menghitung alamat fisik dari pasangan
          segmen dan offset.
    \item[\checkbox] Saya dapat menyebutkan fungsi bus alamat, bus
          data, bus kontrol, dan sinyal \reg{ALE}, \code{M/IO},
          \reg{RD}, \reg{WR}, \reg{READY}, \reg{RESET}.
    \item[\checkbox] Saya dapat menelusuri satu instruksi melalui
          tahapan \eng{fetch}--\eng{decode}--\eng{execute}--\eng{store}.
    \item[\checkbox] Saya dapat menjelaskan mengapa instruksi yang
          mengakses memori lebih lambat daripada instruksi yang
          hanya memakai register.
\end{itemize}

\vspace{0.4em}
Bila masih ada butir yang belum bertanda, ulangi bagian yang
bersesuaian sebelum melanjutkan ke bab berikutnya.
\end{checklistbox}

\begin{infobox}
\textbf{Rujukan bab.} Materi pada bab ini diperkuat oleh
\cite{intel1979} sebagai rujukan utama datasheet dan pembagian
BIU/EU, \cite{stallings2016} untuk gambaran organisasi dan
arsitektur komputer, \cite{courses2103} untuk penjelasan segmentasi
dan alamat fisik, serta \cite{carleton2011} untuk konteks
keluarga x86. Untuk penelusuran lebih dalam mengenai bus dan sinyal
kendali 8086, lihat \cite{intel1979}.
\end{infobox}

\section{Evaluasi}
\label{sec:evaluasi-bab4}

\subsection{Soal Pilihan Ganda}

\begin{enumerate}
    \item Unit yang bertanggung jawab menghitung alamat fisik 20-bit
          dari pasangan segmen:offset pada 8086 adalah \dots
          \begin{enumerate}
              \item \eng{Execution Unit}
              \item \eng{Bus Interface Unit}
              \item \eng{Arithmetic Logic Unit}
              \item register \eng{flag}
          \end{enumerate}

    \item Ukuran \eng{instruction queue} pada mikroprosesor 8086
          adalah \dots
          \begin{enumerate}
              \item 2 byte \item 4 byte
              \item 6 byte \item 16 byte
          \end{enumerate}

    \item Bila \reg{CS} $=$ \hex{1000} dan \reg{IP} $=$ \hex{0025},
          maka instruksi berikutnya diambil dari alamat fisik \dots
          \begin{enumerate}
              \item \hex{10025} \item \hex{10250}
              \item \hex{1002} \item \hex{100025}
          \end{enumerate}

    \item Sinyal yang menandai bahwa pin \reg{AD0}--\reg{AD15}
          sedang membawa alamat sehingga nilainya perlu ditangkap
          \eng{latch} adalah \dots
          \begin{enumerate}
              \item \reg{READY} \item \reg{ALE}
              \item \code{M/IO} \item \reg{RESET}
          \end{enumerate}

    \item Setelah \reg{RESET} aktif, 8086 memulai eksekusi dari
          alamat fisik \dots
          \begin{enumerate}
              \item \hex{00000} \item \hex{10000}
              \item \hex{FFFF0} \item \hex{FFFFF}
          \end{enumerate}
\end{enumerate}

\subsection{Soal Uraian}

\begin{enumerate}
    \item Gambarkan diagram blok 8086 yang memuat EU, BIU, bus
          internal, dan bus eksternal. Sebutkan register yang berada
          di dalam masing-masing unit dan jelaskan mengapa EU tidak
          boleh mengakses bus luar secara langsung.

    \item Jelaskan mekanisme \eng{pipelining} pada 8086 dengan
          bantuan contoh tiga instruksi. Hitung perbedaan jumlah
          siklus antara eksekusi tanpa antrean dan dengan antrean,
          lalu jelaskan kapan keuntungan tersebut hilang.

    \item Sebuah instruksi \lstasm|MOV AX, [BX]| memerlukan waktu
          lebih lama daripada \lstasm|MOV AX, BX|. Uraikan tahapan
          siklus instruksi untuk kedua instruksi tersebut dan
          tunjukkan pada tahap mana perbedaan waktunya muncul.

    \item Diketahui tiga instruksi berturut-turut bertempat pada
          alamat \hex{10000}, \hex{10003}, dan \hex{10005} dengan
          panjang berturut-turut 3 byte, 2 byte, dan 2 byte. Hitung
          nilai \reg{IP} setelah ketiga instruksi selesai bila
          \reg{CS} $=$ \hex{1000} dan \reg{IP} mula-mula \hex{0000}.
          Jelaskan hubungan antara panjang instruksi dan kenaikan
          \reg{IP}.
\end{enumerate}

\subsection{Kisi-kisi Penilaian}

\begin{table}[H]
    \centering
    \caption{Kisi-kisi evaluasi Bab~\ref{chap:arsitektur-8086}.}
    \label{tab:kisi-bab4}
    \small
    \begin{tabular}{@{}p{1.5cm}p{5.0cm}p{2.6cm}c@{}}
        \toprule
        \textbf{Sub-CPMK} & \textbf{Indikator yang diukur} &
        \textbf{Bentuk soal} & \textbf{Bobot} \\
        \midrule
        1.1 & Menjelaskan pembagian tugas BIU dan EU beserta
        register masing-masing & Pilihan ganda 1; Uraian 1 & 30\% \\
        \addlinespace
        1.1 & Menjelaskan \eng{instruction queue} dan
        \eng{pipelining} instruksi & Pilihan ganda 2; Uraian 2 & 25\% \\
        \addlinespace
        1.1 & Menjelaskan bus alamat, data, kontrol dan sinyal
        kendali & Pilihan ganda 4, 5 & 20\% \\
        \addlinespace
        1.1 & Menghitung alamat fisik dan menelusuri siklus
        instruksi & Pilihan ganda 3; Uraian 3, 4 & 25\% \\
        \bottomrule
    \end{tabular}
\end{table}

\begin{checklistbox}
\textbf{Checklist Pencapaian Kompetensi Bab~\ref{chap:arsitektur-8086}}

Seorang mahasiswa dinyatakan kompeten pada Sub-CPMK 1.1 apabila:
\begin{itemize}
    \item[\checked] dapat membedakan arsitektur internal dan
          eksternal serta menyebutkan contoh keduanya;
    \item[\checked] dapat menjelaskan pembagian tugas BIU dan EU
          beserta alasan pemisahan tersebut;
    \item[\checked] dapat menghitung alamat fisik dari pasangan
          segmen dan offset dengan benar;
    \item[\checked] dapat menjelaskan peran \eng{instruction queue}
          enam byte dan pengaruh \eng{pipelining} pada kinerja;
    \item[\checked] dapat menelusuri satu instruksi melalui tahapan
          \eng{fetch}--\eng{decode}--\eng{execute}--\eng{store}
          beserta perubahan \reg{CS}, \reg{IP}, dan antrean;
    \item[\checked] dapat menyebutkan fungsi sinyal \reg{ALE},
          \code{M/IO}, \reg{RD}, \reg{WR}, \reg{READY},
          \reg{RESET}, dan \code{MN/MX}.
\end{itemize}
\end{checklistbox}


\end{document}
